\section{Preliminaries and Problem Definition}\label{sec:preliminaries}
\subsection{Definitions}\label{sec:definitions}

\begin{kartdefinition}[Trajectory]
A trajectory $T_i=\langle p_{i,1},\ldots,p_{i,m_i}\rangle$ is an ordered sequence of spatio-temporal points, where $p_{i,j}=(x_{i,j},y_{i,j},t_{i,j})$.
  Its minimum bounding rectangle, $\operatorname{MBR}(T_i)$, is the smallest axis-aligned rectangle enclosing its locations.
\end{kartdefinition}

\begin{kartdefinition}[Trajectory Query Types]
Let $\mathcal{T}$ be a trajectory dataset, $W_s$ a spatial window, $W_t$ a temporal window, $T_q$ a reference trajectory, $d(\cdot,\cdot)$ a trajectory distance, and $\tau$ a distance threshold.
  The query types are defined below.
\end{kartdefinition}

\noindent\textbf{Spatial range query (SRQ).}
An SRQ returns trajectories with a location in $W_s$.
  Writing $\operatorname{loc}(p)$ for the spatial component of $p$, we have
\begin{equation}\label{eq:srq}
\mathcal{T}_{\mathrm{SRQ}}(W_s)
=\{T_i\in\mathcal{T}\mid \exists p\in T_i:\operatorname{loc}(p)\in W_s\}.
\end{equation}

\noindent\textbf{Spatio-temporal range query (ST-range).}
An ST-range requires a location and its timestamp to satisfy the two windows jointly:
\begin{equation}\label{eq:strange}
\mathcal{T}_{\mathrm{ST}}(W_s,W_t)
=\{T_i\in\mathcal{T}\mid \exists p\in T_i:p\in W_s\times W_t\}.
\end{equation}

\noindent\textbf{Similarity search (SS).}
Given $T_q$ and $\tau$, an SS returns trajectories whose distance to $T_q$ does not exceed $\tau$:
\begin{equation}\label{eq:ss}
\mathcal{T}_{\mathrm{SS}}(T_q,\tau)
=\{T_i\in\mathcal{T}\mid d(T_i,T_q)\le\tau\}.
\end{equation}

\noindent\textbf{Top-$k$ query.}
Given $T_q$ and $1\le k\le|\mathcal{T}|$, a Top-$k$ query returns the $k$ closest trajectories:
\begin{equation}\label{eq:topk}
\mathcal{T}_{\mathrm{Top}\text{-}k}(T_q)
=\arg\min_{\substack{S\subseteq\mathcal{T}\\|S|=k}}
\sum_{T_i\in S}d(T_i,T_q).
\end{equation}

TBD: the reason why KART use SRQ as workload experiments.

\begin{kartdefinition}[Logical Query Plan]
A logical plan $\pi$ is an operator graph that retrieves candidate chunks, refines them against the query, and forms the requested result.
  Its evaluation is written as
\begin{equation}\label{eq:logical-plan}
\llbracket\pi\rrbracket_{\mathcal{D}}(q)
=\operatorname{Post}_{q}\!\left(\operatorname{Refine}_{q}\!\left(\operatorname{Fetch}_{\mathcal{D}}\!\left(\operatorname{Combine}\!\left(\operatorname{Scan}_{a_1},\ldots,\operatorname{Scan}_{a_r}\right)\right)\right)\right),
\end{equation}
where each $\operatorname{Scan}_{a_j}$ uses access path $a_j$, $\operatorname{Combine}$ merges candidate chunks, and $\operatorname{Post}_{q}$ applies any query-specific similarity or Top-$k$ operation.
\end{kartdefinition}

\begin{kartdefinition}[Safe Query Plan]
A logical plan $\pi$ is safe if its access operators retrieve every chunk that may contain a query match.
  Let $\mathcal{C}^{\star}(q,\mathcal{D})$ be those required chunks and $\mathcal{C}_{\pi}(q,\mathcal{D})$ the chunks retrieved by $\pi$. Coverage safety requires
\begin{equation}\label{eq:coverage}
\mathcal{C}^{\star}(q,\mathcal{D})
\subseteq\mathcal{C}_{\pi}(q,\mathcal{D}).
\end{equation}
TBD: in ablation experiments will discuss why without verify is ok
\end{kartdefinition}

\subsection{Query Optimization}\label{sec:optimization}

\noindent\textbf{Access paths and plan alternatives.}
Access paths and execution plans. Access paths map query predicates to key ranges
or equality lookups. An execution plan may use a single index, combine the results of compatible indexes, or
scan the raw data. Kart provides a variety of indexes, including advanced indexes such as TShape and XZ*.

Rule-based optimization enforces a fixed access-path order. 
Conversely, cost-based optimization evaluates estimated latencies, recognizing that while index scans reduce candidate refinement, their scanning and set-operation overhead means the most selective path does not always yield the lowest latency.

\noindent\textbf{Cost model.}
Let $F_{\boldsymbol{\theta}}(\mathbf{x})$ estimate a plan's execution latency from features $\mathbf{x}$ and calibrated coefficients $\boldsymbol{\theta}$.
  The estimate combines index access, set operations, data retrieval, exact filtering, trajectory reconstruction, similarity computation, and Top-$k$ selection:
\begin{equation}\label{eq:cost}
\begin{aligned}
\widehat{L}_{\mathrm{exec}}
={}&\widehat{L}_{\mathrm{index}}+\widehat{L}_{\mathrm{set}}
+\widehat{L}_{\mathrm{fetch}}+\widehat{L}_{\mathrm{exact}}\\
&+\widehat{L}_{\mathrm{reconstruct}}+\widehat{L}_{\mathrm{sim}}
+\widehat{L}_{\mathrm{topk}}.
\end{aligned}
\end{equation}
These metrics are derived from the execution plan, snapshot statistics, and sample estimates. In distributed environments, these metrics also take into account concurrency and the capacity of shared RegionServers.
Coefficient-related terms that are not used are set to 0.
Later, in the “TBD” section, we discuss work on calibrating the coefficients using an offline LLM.

\noindent\textbf{Fast Cost and Final Cost.}
Both estimators use the same model and the feature vector $\mathbf{x}$ from Eq.~\eqref{eq:cost}:
\begin{equation}\label{eq:fast-final}
\begin{aligned}
\fastcost(\pi)&=F_{\boldsymbol{\theta}}(\mathbf{x}(\pi,q,\mathcal{D})),\\
\finalcost(\pi)&=F_{\boldsymbol{\theta}}(\mathbf{x}(\mathcal{K}(\pi),q,\mathcal{D})).
\end{aligned}
\end{equation}
During search, Fast Cost uses the plan, query, and snapshot to estimate ranges and candidate counts. Final Cost uses the compiled access information to refine access and scheduling costs; row and candidate counts remain estimates, so it ranks plans without predicting exact runtime.

\noindent\textbf{Planning objective.}
For bound query $q$ and snapshot $\mathcal{D}$, let $\Pi_{\mathrm{safe}}^{B}(q,\mathcal{D})$ be the safe candidates reached within planning budget $B$.
  Cost-based selection seeks

\begin{equation}\label{eq:objective}
\widehat{\pi}=
\arg\min_{\pi\in\Pi_{\mathrm{safe}}^{B}(q,\mathcal{D})}
\finalcost(\pi).
\end{equation}

Kart uses this cost model for baseline search and constrained rewriting decisions;
Its offline LLM-assisted cost patching and rewriting rules will be detailed later in the TBD section.

\subsection{Problem Definition and Objective}\label{sec:workflow}

Given a natural-language request $u$ and a pinned snapshot $\mathcal{D}$, \KART{} binds $u$ to a query $q$ and selects a safe plan from the candidates reached within budget $B$, minimizing estimated execution cost as in Eq.~\eqref{eq:objective}. Only plans that pass validation are executed, and coverage safety ensures that exact refinement can produce the complete result for $q$.

